%% 
%% Copyright 2007-2025 Elsevier Ltd
%% 
%% This file is part of the 'Elsarticle Bundle'.
%% ---------------------------------------------
%% 
%% It may be distributed under the conditions of the LaTeX Project Public
%% License, either version 1.3 of this license or (at your option) any
%% later version. The latest version of this license is in
%% http://www.latex-project.org/lppl.txt
%% and version 1.3 or later is part of all distributions of LaTeX
%% version 1999/12/01 or later.
%% 
%% The list of all files belonging to the 'Elsarticle Bundle' is
%% given in the file `manifest.txt'.
%% 
%% Template article for Elsevier's document class `elsarticle'
%% with numbered style bibliographic references
%% SP 2008/03/01
%% $Id: elsarticle-template-num.tex 272 2025-01-09 17:36:26Z rishi $
%%
\documentclass[preprint,12pt]{elsarticle}

%% Use the option review to obtain double line spacing
%% \documentclass[authoryear,preprint,review,12pt]{elsarticle}

%% Use the options 1p,twocolumn; 3p; 3p,twocolumn; 5p; or 5p,twocolumn
%% for a journal layout:
%% \documentclass[final,1p,times]{elsarticle}
%% \documentclass[final,1p,times,twocolumn]{elsarticle}
%% \documentclass[final,3p,times]{elsarticle}
%% \documentclass[final,3p,times,twocolumn]{elsarticle}
%% \documentclass[final,5p,times]{elsarticle}
%% \documentclass[final,5p,times,twocolumn]{elsarticle}

%% For including figures, graphicx.sty has been loaded in
%% elsarticle.cls. If you prefer to use the old commands
%% please give \usepackage{epsfig}

%% The amssymb package provides various useful mathematical symbols
\usepackage{amssymb}
%% The amsmath package provides various useful equation environments.
\usepackage{amsmath}
%% The amsthm package provides extended theorem environments
%% \usepackage{amsthm}
\usepackage{xcolor}
%% The lineno packages adds line numbers. Start line numbering with
%% \begin{linenumbers}, end it with \end{linenumbers}. Or switch it on
%% for the whole article with \linenumbers.
\usepackage{lineno}
\usepackage{subcaption}
\usepackage{hyperref}
\usepackage{cleveref}
\usepackage{booktabs}
\usepackage{tabularx}
\usepackage{array}
\usepackage{siunitx}

\journal{Spectrochimica Acta Part A: Molecular and Biomolecular Spectroscopy}

\begin{document}

\begin{frontmatter}

%% Title, authors and addresses

%% use the tnoteref command within \title for footnotes;
%% use the tnotetext command for theassociated footnote;
%% use the fnref command within \author or \affiliation for footnotes;
%% use the fntext command for theassociated footnote;
%% use the corref command within \author for corresponding author footnotes;
%% use the cortext command for theassociated footnote;
%% use the ead command for the email address,
%% and the form \ead[url] for the home page:
%% \title{Title\tnoteref{label1}}
%% \tnotetext[label1]{}
%% \author{Name\corref{cor1}\fnref{label2}}
%% \ead{email address}
%% \ead[url]{home page}
%% \fntext[label2]{}
%% \cortext[cor1]{}
%% \affiliation{organization={},
%%    addressline={},
%%    city={},
%%    postcode={},
%%    state={},
%%    country={}}
%% \fntext[label3]{}

\title{Label-free Raman spectroscopic profiling of circulating extracellular
 vesicles in ischemic stroke: subject-level compositional mapping and the
 limits of individual clinical prediction}

%% use optional labels to link authors explicitly to addresses:
%% \author[label1,label2]{}
%% \affiliation[label1]{organization={},
%%    addressline={},
%%    city={},
%%    postcode={},
%%    state={},
%%    country={}}
%%
%% \affiliation[label2]{organization={},
%%    addressline={},
%%    city={},
%%    postcode={},
%%    state={},
%%    country={}}

\author[1]{[Author 1]\corref{cor1}} %% Author name
\ead{[author1@email.com]}
\cortext[cor1]{Corresponding author}

\author[1,2]{[Author 2]}

%% Author affiliation
\affiliation[1]{organization={[Department, Institution]},
   city={[City]},
   country={[Country]}}
   
\affiliation[2]{organization={[Department, Institution]},
   city={[City]},
   country={[Country]}}

%% Abstract
\begin{abstract}
Extracellular vesicles (EVs) reflect systemic physiological states and hold promise as liquid biopsy analytes. Raman spectroscopy provides a powerful, label-free method for profiling the complex biomolecular cargo of these EVs. However, translating Raman-based EV profiling to clinical settings is frequently hindered by statistical pseudoreplication, which treats multiple pixel-level spectra from the same individual as independent, thereby artificially inflating statistical power. To address this, we developed a subject-level biological principal component analysis (PCA) pipeline and applied it to Raman spectra (fingerprint and high-wavenumber C\textendash H stretching regions) of circulating EVs from 97 ischemic stroke patients and 50 healthy controls. Data were rigorously preprocessed and hierarchically aggregated to yield a single representative spectrum per participant. Subject-level PCA revealed that inter-individual spectral variance is overwhelmingly dominated by lipid compositional differences, explaining up to 80.4\% of the variance. Furthermore, comprehensive intra-disease correlation analyses and cross-validated multivariate regression established that these bulk EV Raman signatures are not predictive of individual cognitive trajectories or clinical endpoints in the stroke cohort. We conclude that bulk circulating EV Raman profiles primarily capture underlying lipid variability, highlighting the necessity for affinity-based target enrichment or single-vesicle analysis in future translational applications.
\end{abstract}

%%Graphical abstract
% \begin{graphicalabstract}
% \includegraphics{grabs}
% \end{graphicalabstract}

% %%Research highlights
% \begin{highlights}
% \item Research highlight 1
% \item Research highlight 2
% \end{highlights}

%% Keywords
\begin{keyword}
Raman spectroscopy \sep extracellular vesicles \sep stroke \sep principal component analysis

%% keywords here, in the form: keyword \sep keyword

%% PACS codes here, in the form: \PACS code \sep code

%% MSC codes here, in the form: \MSC code \sep code
%% or \MSC[2008] code \sep code (2000 is the default)

\end{keyword}

\end{frontmatter}

%% Add \usepackage{lineno} before \begin{document} and uncomment 
%% following line to enable line numbers
%% \linenumbers

%% main text
%%

\section{Results and Discussion}
\label{sec:results}

\subsection{Spectral characteristics of circulating EVs and severity-associated variations}
\label{subsec:spectral_profiles}

\cref{fig:sp_mean_1300} and \cref{fig:sp_mean_2850} display the mean Raman spectra of circulating EVs in the fingerprint and high-wavenumber C\textendash H stretching regions, respectively. Both spectral domains exhibit the classical vibrational profile of a complex, lipid- and protein-rich biological matrix. In the fingerprint region (\cref{fig:sp_mean_1300}), the principal spectral features include the phenylalanine ring breathing mode at $\approx$1003~\si{\per\centi\metre}, the DNA/RNA phosphodiester symmetric stretch near $\approx$1090~\si{\per\centi\metre}, the C\textendash C skeletal stretch at $\approx$1127~\si{\per\centi\metre}, the amide~III envelope (1200--1270~\si{\per\centi\metre}), the CH$_2$ twisting deformation ($\approx$1300~\si{\per\centi\metre}), the CH$_2$/CH$_3$ scissoring/deformation complex ($\approx$1450~\si{\per\centi\metre}), and the amide~I band ($\approx$1657~\si{\per\centi\metre})~\cite{kaminska2022}. In the C\textendash H stretching region (\cref{fig:sp_mean_2850}), the dominant features arise from the symmetric and asymmetric CH$_2$ and CH$_3$ stretching vibrations (2850--2940~\si{\per\centi\metre}) originating from lipid acyl chains and protein methyl groups.

\begin{figure}[htbp]
\centering
\includegraphics[width=\linewidth]{figures/fig_mean_spectra_1300.png}
\caption{Group mean Raman spectra in the fingerprint region (770--1800~\si{\per\centi\metre}). Solid lines represent the group mean; shaded envelopes indicate one standard deviation across all subjects. Key band assignments are annotated.}
\label{fig:sp_mean_1300}
\end{figure}

\begin{figure}[htbp]
\centering
\includegraphics[width=\linewidth]{figures/fig_mean_spectra_2850.png}
\caption{Group mean Raman spectra in the high-wavenumber C\textendash H stretching region (2400--3200~\si{\per\centi\metre}). Solid lines represent the group mean; shaded envelopes indicate one standard deviation.}
\label{fig:sp_mean_2850}
\end{figure}

To determine whether ischemic stroke severity imprints a distinct spectral signature upon circulating EVs, we computed the difference spectra between the severe and mild stroke cohorts (\cref{fig:diff_spectrum}). While the normalised mean intensities of the two groups overlap extensively upon visual inspection, an analysis of the effect sizes (Cohen's $d$) elucidates specific compositional shifts driven by stroke severity. Within the fingerprint region, EVs isolated from patients with severe stroke exhibit relative intensity enhancements in the amide~III region (1200--1300~\si{\per\centi\metre}) and the CH$_2$ scissoring/deformation band ($\approx$1450~\si{\per\centi\metre}). These spectral deviations are indicative of structural alterations in protein secondary conformation and lipid acyl chain packing. Conversely, the mild stroke cohort demonstrates higher relative intensities in the C\textendash C skeletal stretch region (1050--1150~\si{\per\centi\metre}) and the 1500--1600~\si{\per\centi\metre} domain.

In the high-wavenumber region, the spectral deviations are more subtle; however, severe stroke is associated with positive effect sizes proximate to the primary symmetric and asymmetric CH$_2$ stretching modes (2850--2900~\si{\per\centi\metre}). These severity-dependent variations suggest that, although the global spectral architecture remains highly conserved, pronounced clinical severity induces specific alterations in the lipid and protein organisation of the circulating EV pool.

\begin{figure}[htbp]
\centering
\includegraphics[width=\linewidth]{figures/fig_difference_spectrum.png}
\caption{Difference spectrum analysis comparing the severe and mild stroke groups. The upper panels depict the normalised mean Raman intensity for the fingerprint (left) and C\textendash H stretching (right) regions. The lower panels display the signed effect size ($\text{sign}(\Delta) \times |\text{Cohen's } d|$) across the spectral range, delineating specific wavenumbers associated with disease severity.}
\label{fig:diff_spectrum}
\end{figure}

\subsection{Spectral correlations with clinical endpoints and cognitive trajectories}
\label{subsec:correlations}

To further interrogate the relationship between EV biomolecular cargo and clinical outcomes, we mapped the correlation landscape across the spectral features. \cref{fig:corr_heatmap} presents the Spearman correlation heatmaps for the fingerprint and C\textendash H stretching regions against key clinical and demographic variables.

Rather than exhibiting a uniform global alteration, the heatmaps reveal a highly localised correlational architecture. For the cognitive assessments, the baseline and 3-month Montreal Cognitive Assessment (MoCA) scores demonstrate relatively diffuse correlation patterns. Crucially, however, the longitudinal trajectory of cognitive change ($\Delta$MoCA) yields a distinct cluster of positive correlation in the fingerprint region, centred squarely on the amide~III band (1214--1364~\si{\per\centi\metre}). This observation parallels the severity-associated increases identified in the difference spectra, implying that protein conformational changes within the EV pool may co-vary with the biological mechanisms driving cognitive decline.

Furthermore, the Repeatable Battery for the Assessment of Neuropsychological Status (RBANS) total score exhibits discrete correlational clusters. Positive associations emerge near 1510~\si{\per\centi\metre}---potentially overlapping with carotenoid C=C stretching---and within the 2895--2998~\si{\per\centi\metre} domain in the high-wavenumber region, which captures the CH$_3$ stretching commonly associated with protein methyl groups.

Importantly, baseline demographics---particularly age---profoundly sculpt the spectral landscape. Age exhibits pronounced correlation clusters, notably a robust positive band extending from 1510 to 1651~\si{\per\centi\metre}, interspersed with alternating negative regions. This pervasive influence emphasises that systemic metabolic and demographic factors act as major drivers of EV compositional heterogeneity.

\begin{figure}[htbp]
\centering
\includegraphics[width=\linewidth]{figures/fig_clinical_corr_heatmap.png}
\caption{Spearman correlation heatmaps delineating the relationship between subject-level Raman spectral features and clinical endpoints. The upper panel illustrates the fingerprint region (733--1789~\si{\per\centi\metre}), while the lower panel displays the C\textendash H stretching region (2450--3198~\si{\per\centi\metre}). Colour intensity corresponds to the magnitude and direction of the Spearman rank correlation coefficient ($\rho$).}
\label{fig:corr_heatmap}
\end{figure}

\subsection{Biochemical context and implications for clinical translation}
\label{subsec:biochem}

Direct inspection of the spectral data establishes that the circulating EV pool harbours severity-associated and cognition-associated biochemical signatures, most notably within the amide~III, CH$_2$/CH$_3$ stretching, and C=C regions. Nevertheless, the overwhelming variance observed in bulk EV preparations remains dictated by systemic lipid heterogeneity. This conclusion is strongly supported by the dominant CH$_2$ features and the pervasive influence of patient age on the correlation profiles.

Although ischemic stroke precipitates the release of brain-derived EVs (BDEVs) carrying neuronal and glial cargo, these particles constitute only a minor fraction of the total circulating pool, which is heavily dominated by platelet-, erythrocyte-, and endothelial-derived vesicles~\cite{liu2025}. The severity-dependent intensity enhancements in structural bands ($\approx$1250, $\approx$1450, and 2850--2900~\si{\per\centi\metre}) and their correlation with $\Delta$MoCA likely reflect a generalised systemic inflammatory and metabolic response to cerebrovascular injury, rather than an isolated neuro-specific signal.

While these spectral correlations yield valuable biological insights, the relatively modest magnitude of these shifts underscores the fundamental limitations of employing bulk EV Raman spectroscopy for robust individual-level clinical prediction. Advanced multivariate predictive modelling applied previously to this dataset confirmed that these subtle features fail to overcome the systemic noise floor, yielding predictions that do not surpass cohort-mean estimates. Consequently, future translational efforts must focus on surmounting this biological dilution. Realising reliable clinical predictions will unequivocally necessitate either the affinity-based enrichment of specific EV subpopulations (e.g., utilising neural surface markers) prior to Raman acquisition~\cite{liu2025}, or the deployment of single-vesicle spectroscopic techniques~\cite{jensen2024, koster2021} to effectively isolate the target pathological signatures from the systemic background.

\section{Conclusions}
\label{sec:conclusions}

This study has presented a rigorous, subject-level Raman spectroscopic analysis of circulating EVs in 97 ischemic stroke patients and 50 healthy controls. The following conclusions are drawn from the data:

\begin{enumerate}
  \item A multi-stage preprocessing pipeline incorporating hierarchical median aggregation effectively eliminates pseudoreplication, enabling the robust characterization of inter-individual EV variability.
  \item Difference spectrum analysis reveals that severe ischemic stroke is associated with distinct compositional shifts in circulating EVs, notably an enhancement in the relative intensity of the amide~III region (1200--1300~\si{\per\centi\metre}) and the CH$_2$ scissoring/deformation bands ($\approx$1450~\si{\per\centi\metre}). This suggests alterations in protein secondary structure and lipid chain packing.
  \item Spectral correlation heatmaps demonstrate that the trajectory of cognitive change ($\Delta$MoCA) mirrors these severity-associated shifts, localizing primarily to the amide~III band. Age and other baseline demographics also strongly pattern the spectral landscape, underscoring the profound systemic metabolic influence on EV composition.
  \item Despite these targeted insights, the overwhelming inter-individual variance remains driven by systemic lipid heterogeneity. The severity-dependent signals likely reflect a generalized systemic inflammatory and metabolic response rather than an isolated neuro-specific signature.
  \item Future clinical applications of EV Raman spectroscopy in cerebrovascular disease must prioritize the affinity-based isolation of brain-derived EV subpopulations or the implementation of single-vesicle Raman analysis to overcome the profound signal dilution imposed by the dominant lipid background of non-cerebral circulating EVs.
\end{enumerate}

\section*{Data Availability}

The processed spectral data matrices and analysis scripts supporting the findings of this study are available from the corresponding author on reasonable request. Raw Raman data are subject to local data governance restrictions.

\section*{Ethics Declaration}

The study protocol received approval from the relevant Institutional Review Board ([IRB reference]). Written informed consent was obtained from all participants or their legal representatives.

\section*{Author Contributions}

[CRediT taxonomy statement to be completed by authors]

\section*{Declaration of Competing Interests}

The authors declare no competing interests.

\section*{Funding}

[Funding statement to be completed by authors]

\section*{Acknowledgements}

The authors thank [to be completed by authors].

\bibliographystyle{elsarticle-num} 
\bibliography{cas-refs}

\end{document}


